\documentclass[11pt]{article}
\usepackage[margin=.87in]{geometry}
\usepackage[T1]{fontenc}
\usepackage{lmodern,microtype,amsmath,amssymb,graphicx,booktabs,tabularx}
\usepackage{xcolor}
\usepackage[colorlinks=true,linkcolor=blue!40!black,urlcolor=blue!50!black]{hyperref}
\setlength{\parindent}{0pt}
\setlength{\parskip}{.4em}
\setlength{\emergencystretch}{1em}
\graphicspath{{images/}}
\title{An effective contact for stock redistribution and a finite clock-energy supply}
\author{Research calculation for the Entropic Scalar EFT project}
\date{23 October 2026}
\begin{document}
\maketitle
\begin{abstract}
A selected effective contact connects the positive depletion completion to
orbital redistribution. A local radial-velocity kernel preserves particle
number and angular momentum and relaxes toward a supplied positive orbital
distribution. Its energy exchange separates mechanical work from bath heat.
Coupling the work to the source clock with power proportional to the square
of the clock speed preserves the regular positive-current equation. A
Hermitian joint transition provides a finite-energy realization: orbital,
work-register and bath energies change together, and a symmetric amplitude
suppresses transitions at the clock floor. A finite-dimensional Hamiltonian
checks that construction. Particle tests with clock feedback reach the
requested distribution when sufficiently supplied and remain far from it
when the clock supply is exhausted. The target distribution, physical rate,
bath spectrum and identification of the work register remain effective inputs.
Spatial scalar tests converge during smooth early evolution; sharp later
fronts retain a resolved-energy deficit. A selected thermal fluid extension
then supplies a conservative continuation with positive energy and causal
local sound speed. Its warm contact gate reduces exactly to the cold rule.
The thermal equation of state and its relation to the anisotropic kinetic
stock remain additional assumptions. The construction is not a new galaxy
fit or a derivation of the full covariant multistream theory.
\end{abstract}

\section{What has been connected}
The preceding research note constructed a positive low-density completion
of the capacity-clock action. This calculation adds a specified interaction,
its work and heat ledger, a finite-energy Hamiltonian realization, and tests
of actual orbital transport. These are effective constructions. The
microstructure has not selected their target or matrix elements.

The other chat's assessment correctly distinguishes source-clock slowing
from ordinary proper time. Throughout this note, $v=d\phi/d\tau$ is the
scalar's speed along its regular flow. It is not the speed of atomic clocks,
biological processes or all physics in an underdense region.

\begin{table}[ht]
\small
\begin{tabularx}{\textwidth}{@{}p{.23\textwidth}XX@{}}
\toprule
Part & Supplied here & Still required \\
\midrule
Exhaustion and work & A regular gated source and finite supply bound &
Its microscopic selection and spatial stress completion \\
Redistribution & Local velocity transitions with a prescribed stationary
orbital distribution & Selection and nonlinear self-consistency of that distribution \\
Reaction and heat & Joint energy shifts and opposite radial impulses &
Dynamics of the bath stress and finite carrier \\
Hamiltonian contact & A Hermitian energy-conserving work-register construction &
Identification with the native clock and a bath spectrum giving the rate \\
Spatial boundary & A conservative thermal fluid option with causal sound
speed & Selection of its stress closure and realistic void predictions \\
\bottomrule
\end{tabularx}
\caption{The new construction supplies a usable conditional interaction.
Its physical inputs and unresolved spatial completion remain explicit.}
\end{table}

\section{The work coupling required near exhaustion}
Use $(-,+,+,+)$ signature and $X=-\nabla\phi\cdot\nabla\phi/2$. The selected
matter action from the preceding note has
\begin{equation}
 P=\lambda X-e(\lambda,\phi)-V(\phi),\qquad
 e=\frac{\lambda}{2}+\frac{2[L-\lambda]_+^3}{3L^2},\qquad
 L=3(K-V)>0,\quad\lambda\ge L/2.
\end{equation}
On its regular timelike pressure branch put $v=\sqrt{2X}$,
$j=\lambda v$, $r=\lambda/L$ and $\theta=\nabla_\mu u^\mu$.
The existing scalar equation is
\begin{equation}
 \dot j+\theta j=-V_\phi A(r)v^2,\qquad
 A(r)=\frac{1+2r-2r^2}{3-2r}\in[3/4,1].
\end{equation}
The dust branch has $v=1$ and source $V_\phi$. It enters the pressure branch
at $j=L>0$. The local energy density has lower bound $K$.

Suppose a contact takes work power $\mathcal P$ per physical volume from
the clock sector, measured in this frame. Its source in the current equation
is $-\mathcal P/v$. A sufficient regularity condition is
\begin{equation}\label{eq:gate}
 \mathcal P=v^2 B,\qquad B\ \hbox{finite and continuous near rest}.
\end{equation}
Then
\begin{equation}\label{eq:current}
 \dot j=-\left[\theta+\frac{V_\phi A(r)j}{\lambda^2}
                         +\frac{B}{\lambda}\right]j.
\end{equation}
The fractional rate remains regular when $K-V$ stays bounded away from
zero. The integrating-factor positivity argument continues to apply.
$B$ can have either sign: the interaction may return work to the clock.
An arbitrary finite power independent of $v$ would reintroduce a singular
source at rest.

This condition selects a class, not one unique interaction. Powers vanishing
faster than $v^2$ also work. A source vanishing more slowly may admit a
different stopping prescription, but does not have the same bounded
fractional-rate argument. On the spacelike branch the proposed contact is
switched off; extension through nonzero null gradients remains part of the
interface problem.

The scalar equation with generalized force is
\begin{equation}
 \nabla_\mu(\lambda\nabla^\mu\phi)-e_\phi-V_\phi=Bv.
\end{equation}
Its stress divergence is therefore
\begin{equation}\label{eq:clockforce}
 \nabla_\mu T_\phi^{\mu\nu}=Bv\nabla^\nu\phi
                         =-\mathcal P u^\nu.
\end{equation}
If $I^\nu$ is the kinetic contact's four-momentum gain, the assignments
\begin{equation}\label{eq:ledger}
 \nabla T_{\rm kin}=I^\nu,\qquad
 \nabla T_{\rm bath}=\mathcal P u^\nu-I^\nu
\end{equation}
close total stress conservation together with~\eqref{eq:clockforce}.
The bath carries the opposite spatial impulse as well as heat. Specifying
these exchange vectors does not determine its constitutive stress or prove
a full covariant initial-value problem.

\subsection{A finite activity budget}
For a homogeneous cell at constant $V$, write its clock Hamiltonian density
as $R(j)$. Legendre duality gives $R'(j)=v$. Let an internal two-state
population have upper-state fraction $w$, target $p$, and rate
\begin{equation}
 \dot w=\nu v^2(p-w).
\end{equation}
If each net upward transition draws work $\mathcal W>0$, then
\begin{equation}
 R(j)+\mathcal W w=R(j_0)+\mathcal W w_0.
\end{equation}
Consequently
\begin{equation}\label{eq:budget}
 w\le w_0+\frac{R(j_0)-K}{\mathcal W}.
\end{equation}
Introduce contact activity $s=\int v^2dt$. The exact population is
$w(s)=p+(w_0-p)e^{-\nu s}$. If the available energy
$D=R(j_0)-K$ is smaller than $\mathcal W(p-w_0)$, then
\begin{equation}
 s_\infty=-\frac1\nu\ln\left[1-\frac{D}{\mathcal W(p-w_0)}\right],
 \qquad w_\infty=w_0+D/\mathcal W<p.
\end{equation}
The current approaches zero and the target is not reached. The mechanism
does not turn the vacuum-like floor into an unlimited source of work.

The test used $K=1$, $p=0.8$, $w_0=0.1$, a unit excitation gap and bath
temperature $\Theta=0.05$ in dimensionless energy units. Detailed balance
below gives $\mathcal W=1+\Theta\ln4=1.0693147$; the required work is
$0.7485203$. With $D=0.25$, the population ends at $0.3337946$, agreeing
with~\eqref{eq:budget}. With $D=1$, it reaches $0.8$. The maximum absolute
energy-ledger error was $2.5\times10^{-11}$. No exhausted state was clipped.

\begin{figure}[ht]
\centering\includegraphics[width=\textwidth]{closed_cell.png}
\caption{Deterministic finite-supply test. The population relaxes in contact
activity rather than ordinary time; activity becomes finite when work is
exhausted. These are dimensionless internal-state tests.}
\end{figure}

\section{A local transition that redistributes orbital stock}
In a fixed spherical potential let
\begin{equation}
 E=\frac{p_r^2}{2\mu}+\frac{\ell^2}{2\mu r^2}+\mu\Phi(r),
\end{equation}
where $\mu$ is a particle or reduced mass. Retain the full angular momentum
vector $\boldsymbol\ell$, including its orbital plane. Supply a positive
integrable stationary radial phase density $h(E,\ell)$, with the same radial
shape for orientations of the same angular-momentum magnitude. At fixed
$r,\boldsymbol\ell$, define
\begin{equation}
 Z(r,\ell)=\int h(E(r,p_r),\ell)\,dp_r,
 \qquad n(r,\boldsymbol\ell)=\int f(r,p_r,\boldsymbol\ell)\,dp_r.
\end{equation}
A selected radial-velocity contact is
\begin{align}
 k(p_r\to p_r')&=\nu v^2\frac{h(E',\ell)}{Z},\label{eq:rate}\\
 C[f]&=\nu v^2\left[\frac{h}{Z}n-f\right].\label{eq:collision}
\end{align}
The rate $\nu$ is independent of the incoming and outgoing radial momenta
within this cell. For the central operator construction it may vary with
radius and angular-momentum magnitude, with no imposed orientation dependence.
Equation~\eqref{eq:collision} is the small-work-per-event limit; finite
clock energies require the symmetric replacement in the next section.

This is a local velocity change. It does not move a particle instantaneously
between radii or impose a new density profile by hand. Hamiltonian streaming
between interactions supplies the spatial transport:
\begin{equation}\label{eq:kinetic}
 \partial_t f+\{f,E\}=C[f].
\end{equation}
The contact preserves positivity, local number and every retained angular
bin. It obeys detailed balance with respect to $h$ when the clock is held
fixed. Since $h$ depends on orbital invariants, Hamiltonian streaming also
leaves it stationary. Relative entropy to that supplied target decreases
under the contact. On a connected accessible domain with nonzero contact,
the joint stationary family is $f=a(\boldsymbol\ell)h$, with each
$a$ fixed by its retained mass. This identifies a transition realizing a
specified equilibrium; it does not derive the equilibrium itself.

In a self-gravitating implementation, the density of this same $f$ sources
the stock potential once. A local velocity reset leaves that density
unchanged at the instant of collision; its later motion changes the density.
Adding a separate full RAR halo would count the excess twice. The stationary
distribution must be consistent with its own potential before the fixed-target
argument can be used. Nonlinear self-gravitating stability does not follow
from the fixed-potential entropy statement.

\subsection{How this uses the existing orbital construction}
The saved positive-orbit witness specifies mass per radial energy interval
at fixed angular momentum. The phase density needed here is related by
\begin{equation}
 h(E,\ell)\ \propto\
 \frac{w_*(E,\ell)}{\mathcal D(E,\ell)},
 \qquad
 \mathcal D(E,\ell)=\int dr\,dp_r\,
       \delta[E-H(r,p_r,\ell)].
\end{equation}
With the corresponding fixed-angular normalization, $\mathcal D$ is the
radial period. Bin widths must be included when the supplied $w_*$ are
discrete masses. A positive interpolation and treatment of singular circular
orbits must be specified. The existing feasible endpoint therefore provides
input for this contact without replacing orbital motion by radial density
diffusion. Its interpolation and a self-consistent evolving halo have not
been implemented in this calculation.

\subsection{Momentum and angular reaction}
A radial relative-coordinate impulse has
$\mathbf p' =\mathbf p+\alpha\mathbf r$ and thus
$\mathbf r\times\mathbf p'=\mathbf r\times\mathbf p$ exactly.
For a finite carrier, opposite particle and carrier impulses conserve total
linear momentum. The reduced-mass kinetic change accounts for their common
mechanical work. Stochastic radial impulses preserve angular momentum
path by path; noise does not merely cancel it on average.

The same property can be stated as a scalar radial operator commuting with
all angular generators. Gravitational torques in a nonspherical potential
can still change an orbit's plane. Carrier assignments and their reaction
stresses must be supplied in a multicarrier system.

The numerical test uses Newtonian orbital motion. A relativistic kinetic
version can bound streaming speed by replacing the energy with its
relativistic dispersion; locality of the collision then avoids the
instantaneous spatial propagation of a diffusion equation. That extension,
and its gravity, were not simulated here.

\section{Work, heat, and an explicit Hamiltonian contact}
\subsection{A nonthermal target requires work}
Let $\Theta=k_BT$ in the same energy units as $E$. The rate ratio is
\begin{equation}
 \ln\frac{k(i\to j)}{k(j\to i)}=\ln\frac{h_j}{h_i}.
\end{equation}
Using local detailed balance\cite{Seifert}, the bath receives heat
\begin{equation}\label{eq:heat}
 Q_{ij}=\Theta\ln(h_j/h_i),\qquad
 W_{ij}=E_j-E_i+Q_{ij}.
\end{equation}
The system gains mechanical energy $W_{ij}-Q_{ij}$.
Without work, the contact requires $h\propto e^{-E/\Theta}$ within each
connected angular bin. A general RAR-compatible orbital target need not
have that form. Calling its contact a passive thermal bath would conceal
the work term.

At a fixed target define
\begin{equation}
 G(E,\ell)=E+\Theta\ln h(E,\ell).
\end{equation}
An arbitrary fixed normalization of $h$ only adds a constant to $G$.
Then $W_{ij}=G_j-G_i$. A work register with energy $U$ changes to
$U'=U-W_{ij}$; $U+G$ is invariant. For many particles sharing the register,
the corresponding invariant is $U+\sum_aG_a$ in total-energy units.
Hamiltonian streaming at fixed spherical potential also preserves each
$G_a$. A moving target or evolving potential requires additional work
accounting and is not covered by this invariant.

\subsection{A symmetric interaction with a lower energy boundary}
Give the effective work register a spectrum $U\ge K$ and let $v(U)$ be
the homogeneous source-clock speed obtained from the selected action.
This identification is an additional effective hypothesis about the
semiclassical register; the native scalar spectrum has not been derived.
Define $a(U)=\sqrt{v(U)}$. For radial system transition $A_{ji}$, work
translation $T_{-W_{ij}}$ and bath operator $B_{Q_{ij}}$ that changes bath
energy by $Q_{ij}$, consider
\begin{equation}\label{eq:Hint}
 H_{\rm ex}=\sum_{ij}g_{ij}\,
 A_{ji}\otimes\bigl[a(\hat U)T_{-W_{ij}}a(\hat U)\bigr]
 \otimes B_{Q_{ij}}+\text{h.c.}
\end{equation}
Translations outside the allowed register spectrum have zero matrix
elements. Angular-degenerate radial transitions are summed identically
over magnetic labels, so the contact commutes with each angular generator.
This construction uses a common scalar register and central radial coupling.
An arbitrary angular variation of the clock or bath is not covered by this
angular selection rule.

Each allowed term changes free energy by
\begin{equation}
 (E_j-E_i)-W_{ij}+Q_{ij}=0.
\end{equation}
Hence the exchange commutes with the joint free Hamiltonian when the stated
energy-resolved operators exist. The forward and reverse register matrix
elements are conjugates. Their squared clock factor is
\begin{equation}\label{eq:edgegate}
 v(U)v(U-W_{ij}),
\end{equation}
which is symmetric between the two endpoint states and vanishes at the
floor. For an individual event small compared with the available register
energy, this reduces to $v(U)^2$ and gives~\eqref{eq:gate}.

This symmetric factor matters. Multiplying only the forward rate by the
initial $v^2$ does not preserve the joint forward/reverse ratio when an
event appreciably changes the register. The earlier one-sided-gate tests
are retained as mean-field diagnostics. The final finite-event test uses
equation~\eqref{eq:edgegate} and rejects only energetically unavailable
proposals; it never clips a negative energy after a transition.

Energy conservation alone does not guarantee that an interaction preserves
the free Hamiltonian's lower bound. To enforce that stronger property,
complete each exchange edge between degenerate joint states $|i\rangle$
and $|f\rangle$ to the positive form
\begin{equation}\label{eq:positiveedge}
 H_c=\sum_e c_e(|i\rangle-|f\rangle)(\langle i|-\langle f|),
 \qquad c_e=|g_e|\sqrt{v(U_i)v(U_f)}\ge0.
\end{equation}
Real edge phases are chosen here; phases can also be included in the
rank-one vector. The diagonal terms are part of the selected interaction.
They provide the least-trace positive completion of a fixed two-state
off-diagonal exchange: a positive block with off-diagonal magnitude $c$
has diagonal entries satisfying $d_1+d_2\ge2c$, attained by $d_1=d_2=c$.
Thus $H_0+H_c\ge H_0$ as quadratic forms, while the free-energy and angular
commutators still vanish. This preserves the free lower bound; it does not
assert that gravitational binding energy has a universal local lower bound.
The additional diagonal terms must be included in any finite-coupling
frequency matching. The rate construction is a weak-coupling effective
description, not an all-orders derivation that can omit these shifts.

Under a suitable weak-coupling reduction, a thermal bath can supply the
ratio $\exp(Q_{ij}/\Theta)=h_j/h_i$. Its spectrum and the matrix elements
$g_{ij}$ determine the absolute rates. At fixed total $U+\sum G_a$, the
ideal joint equilibrium is the supplied orbital weight restricted to
configurations with sufficient remaining register energy, subject to the
chosen energy-normalized register measure. The finite supply therefore
changes the reachable target near exhaustion. It does not leave a second
unlimited controller hidden outside the work ledger.

The construction establishes a possible interaction, not its unique
physical selection. The energy basis and matrix elements, the required
bath correlations, and validity of a Markov approximation remain to be
matched. In particular, the two-pole response in the main paper is not
derived by writing~\eqref{eq:Hint}. Bath KMS symmetry alone also does not
establish the equilibrium fluctuation--dissipation formula for a driven
stock observable. Quantum noise must be checked in this coupled model.

\subsection{A finite Hamiltonian witness}
The positive completion in~\eqref{eq:positiveedge} is used in the final
Hamiltonian witness. The original exchange-only witness is retained to
distinguish conservation from the additional positivity requirement.
A 72-state witness uses two radial states, six work-register levels, two
bath levels and three magnetic labels. A transition raises orbital energy
by $2$, lowers work energy by $3$, and raises bath energy by $1$. The
floor is $U=1$. Hermiticity, the commutator with joint free energy, all
three angular commutators and the coupling out of the floor were zero
numerically. Unitary evolution conserved norm to $6.7\times10^{-16}$ and
free energy to $4.5\times10^{-15}$. The minimum eigenvalue of the full
Hamiltonian remained $1$, the free ground energy; its interaction was
positive within eigensolver precision ($10^{-16}$).
This small closed bath displays coherent exchange. It is not a derivation
of dissipative relaxation or a measurement of a physical rate.

\section{Orbital and coupled tests}
\subsection{Local velocity contact in a fixed potential}
The fixed spherical test uses
$E=p_r^2/2+r^2/2+\ell^2/(2r^2)$, $\ell=0.8$, and
\begin{equation}
 h(E)=0.7e^{-E/0.25}+0.3e^{-E/1}.
\end{equation}
The two numbers in the exponential define the chosen orbital distribution;
they are not both temperatures of the bath. The conditional radial-velocity
distribution is sampled exactly as a position-dependent Gaussian mixture.
Forty thousand particles underwent Hamiltonian streaming and local resets
at common rate $\nu=2$ in dimensionless units.

At steps $0.01$ and $0.005$, the radial cumulative-distribution discrepancies
from the supplied stationary target were $0.00458$ and $0.00340$.
Different steps have different random histories; these are sampling checks,
not deterministic convergence rates. The Hamiltonian integration error
decreased from $2.08\times10^{-4}$ to $5.43\times10^{-5}$. The eventwise
work-minus-heat identity held to $3.6\times10^{-15}$.

\begin{figure}[ht]
\centering\includegraphics[width=\textwidth]{local_contact.png}
\caption{Actual orbital motion plus local radial-velocity transitions.
The declared stationary target is recovered in a fixed potential. This
example is not an RAR fit or a self-consistent halo.}
\end{figure}

\subsection{Finite clock supply with the symmetric contact}
The final coupled test uses 12,000 particles, the same orbital target, a
shared clock-energy register, and the symmetric factor~\eqref{eq:edgegate}.
The bath temperature is $10^{-6}$ in orbital energy units, so thermal heat
is small compared with mechanical work. This is a dimensionless small-
temperature test; it is not a dimensional identification with the horizon
temperature. Earlier higher-temperature tests are also retained.

At the finer step $0.005$, initial work supplies $0.15$ and $1.2$ give:
\begin{center}
\begin{tabular}{@{}lrr@{}}
\toprule
Quantity & Limited supply & Sufficient supply \\
\midrule
Final energy above floor & $5.04\times10^{-5}$ & $0.2292$ \\
Final source-clock speed & $0.00820$ & $0.5035$ \\
Radial CDF discrepancy & $0.5029$ & $0.00758$ \\
Work transferred & $0.14995$ & $0.97081$ \\
Heat to bath & $-1.78\times10^{-7}$ & $-1.11\times10^{-6}$ \\
Maximum energy-ledger error & $2.20\times10^{-6}$ & $1.02\times10^{-5}$ \\
\bottomrule
\end{tabular}
\end{center}
The negative heat entries mean the thermostat supplies those small amounts
of energy. Its energy change is included in the ledger; a finite thermal
reservoir and its spatial stress were not evolved. The lower step reduced
the mechanical integration error in both supply cases.

\begin{figure}[ht]
\centering\includegraphics[width=\textwidth]{hermitian_contact.png}
\caption{The symmetric finite-event contact. The low-supply population
remains far from its requested equilibrium; the well-supplied population
approaches it. No negative clock energy is generated and then repaired.}
\end{figure}

The clock and orbital energies are actually coupled in this test. The
model includes excitation energy once and does not add a second rest-mass
halo. It is nevertheless a coarse energy-level construction in a frozen
potential. Identifying its register, kinetic inertia and bath stress as
parts of one covariant capacity medium remains an additional derivation.
Finite stochastic exchange can return energy to the register; the
near-floor speed need not decline monotonically along every realization.

\section{Spatial boundaries and a conservative thermal continuation}
For constant $V$ in a flat one-dimensional patch, put
$s=\partial_x\phi$ and $J=\lambda\partial_t\phi$. The Hamiltonian equations are
\begin{equation}
 \partial_tJ=\partial_x(\lambda s),\qquad
 \partial_ts=\partial_xv,
\end{equation}
and the smooth energy balance is
\begin{equation}
 \partial_t\mathcal H=\partial_x(Js).
\end{equation}
The constitutive inversion uses the dust, pressure and canonical spacelike
branches from the earlier note. Regular scalar characteristics are those
of a $P(X,\phi)$ fluid\cite{Garriga}. Unequal clock speeds do not form a
stationary planar boundary in this setting: continuity of the second flux
requires equal $v$ for a stationary interface. The interface must evolve.

A compact smooth transition joined $J=0.2$ to $J=3.6$ with initial $s=0$,
$L=3$ and $K=1$. A conservative finite-volume method evolved both interface
sides without imposing positive-current clipping. During the early interval
$t=0.1$, the fractional total-energy errors were
$7.99\times10^{-5}$, $1.96\times10^{-5}$ and $4.88\times10^{-6}$ at
256, 512 and 1024 cells. This is approximately second-order convergence.

By $t=5$, compressive fronts were sharp. At 256, 512, 1024 and 2048 cells,
the resolved scalar energy deficits were respectively $0.607\%$, $0.440\%$,
$0.355\%$ and $0.319\%$. They have not converged to zero over these grids.
The minimum current remained $0.2$, the minimum $X$ remained positive,
and the energy floor held in every tested cell. Mean current was conserved
to $9\times10^{-16}$. Small numerical tails beyond the unit-speed cone
decreased with refinement.

\begin{figure}[ht]
\centering\includegraphics[width=\textwidth]{interface.png}
\caption{The pressure difference drives spatial motion and later compression.
The scalar-only solver preserves the tested positive states but loses
resolved energy at sharp fronts. The plot does not establish a conservative
shock or multistream completion.}
\end{figure}

This is a useful limitation rather than an accepted nonlinear solution.
The numerical viscosity can remove resolved energy once gradients become
sharp. A physical continuation must specify whether that energy enters
unresolved velocity dispersion, heat or another degree of freedom, and
include its stress and feedback. The work-and-bath ledger above identifies
where such energy and momentum would be recorded. The selected thermal
continuation below provides one concrete option. Realistic compensated
voids, curved geometry and multistream formation remain untested.

\subsection{A warm extension of the same cold energy}
Let $n$ denote proper current density and let $R_0(n,\phi)$ be the cold
Hamiltonian obtained from the previous scalar action. Its derivative is
$\mu_0=\partial R_0/\partial n=v_0$, with
$n=\lambda_0v_0$, and $p_0=n\mu_0-R_0$. Introduce a nonnegative advected
coefficient $z$ and choose
\begin{equation}\label{eq:warm}
 \rho=R_0(n,\phi)+e_{\rm th},\qquad
 e_{\rm th}=z n^\gamma,\qquad 1<\gamma\le2.
\end{equation}
Then
\begin{equation}
 p=p_0+(\gamma-1)e_{\rm th},\qquad
 \mu=\left.\frac{\partial\rho}{\partial n}\right|_{z,\phi}
     =\mu_0+\frac{\gamma e_{\rm th}}n,
 \qquad \rho+p=n\mu.
\end{equation}
The new test chooses $\gamma=4/3$, a radiation-like isotropic stress.
The equation of state is an effective choice, not a derived value of a new
microscopic parameter. Thermal energy is determined by the conserved
energy equation from initially cold data; no thermal-density amplitude is
fitted to the interface.

Writing $h_0=n\mu_0$, $h_T=\gamma e_{\rm th}$ and denoting the cold
adiabatic sound speed by $c_0$, the warm sound speed is
\begin{equation}\label{eq:warmsound}
 c_s^2=\frac{h_0c_0^2+h_T(\gamma-1)}{h_0+h_T}.
\end{equation}
It is an enthalpy-weighted average of two numbers in $[0,1]$.
The energy floor also persists: $\rho\ge K$, and for any observer with
Lorentz factor $W$, $T_{00}=\rho+(\rho+p)(W^2-1)\ge K$.
For bounded $z$, the limit $n\to0$ gives
\begin{equation}
 \rho\to K,\qquad p\to-K,\qquad \mu\to0.
\end{equation}
Thus the warm extension retains the vacuum-like resting state. On the
exactly cold, irrotational branch it reduces to the previous action.
Additional entropy modes are possible; inheriting the old linear fit
requires the same cold adiabatic initial sector, with no separately seeded
thermal component.

A local ideal-fluid action can use a current $J^\mu=nu^\mu$ and advected
potentials, in the standard variational fluid framework\cite{Brown}:
\begin{equation}\label{eq:fluidaction}
 S=\int\sqrt{-g}\left[-\rho(n,z,\phi)
 +J^\mu\left(\partial_\mu\phi+
       \sum_A\beta_A\partial_\mu\alpha_A+\zeta\partial_\mu z\right)\right].
\end{equation}
Variation transports $z$ and $\alpha_A$ in smooth adiabatic flow, gives
$u^\mu\partial_\mu\phi=\mu$ and
$\nabla_\mu J^\mu=-\rho_\phi$, and supplies
$T^{\mu\nu}=(\rho+p)u^\mu u^\nu+pg^{\mu\nu}$.
The explicit $\phi$ dependence is inherited from the cold potential.
Entropy-producing shocks require the conservative weak equations and their
entropy condition; they are not smooth extrema of this ideal action.
This is a fluid stress completion with possible vorticity, not a complete
collisionless multistream distribution.

\subsection{The safe gate becomes an enthalpy gate}
The warm clock speed $\mu$ no longer equals $v_0$. For $\gamma<2$,
using $\mathcal P\propto\mu^2$ alone need not keep the fractional current
drain bounded near $n=0$. A bounded choice that retains the cold rule is
\begin{equation}\label{eq:warmgate}
 g=\frac{n\mu}{\lambda_0+h_T},\qquad
 \mathcal P=gB,\qquad 0\le g\le1.
\end{equation}
Indeed $n\mu=h_0+h_T$ and $h_0=\lambda_0v_0^2\le\lambda_0$.
When $h_T=0$, $g=v_0^2$ exactly. For positive current,
\begin{equation}
 \frac{\mathcal P}{\mu}=\frac{nB}{\lambda_0+h_T},
 \qquad
 \dot n=-\left[\theta+
 \frac{V_\phi A\,n}{\lambda_0^2}+\frac{B}{\lambda_0+h_T}\right]n.
\end{equation}
The pressure-branch cold source is unchanged at fixed $n,z$ because the
thermal term has no explicit $\phi$ dependence. For bounded coefficients
and $K-V>0$, the fractional current equation remains regular. In the
Hamiltonian contact, replace $a=v_0^{1/2}$ by $a=g^{1/4}$; the symmetric
rate factor becomes $\sqrt{g_i g_f}$ and its small-event limit is $g$.
This is a specified warm extension of the coupling, with no extra fitted
amplitude. It does not make this rational gate unique among admissible
positive gates.

For a vortical or nonisentropic fluid, the full potential gradients need
not be parallel to $u^\mu$. The momentum delivered by generalized forces
must therefore be retained in the bath ledger; its energy projection
alone is insufficient. The warm gate establishes the regular energy/current
balance, not a complete bath-force action.

\subsection{Conservative compression test}
The same interface was evolved using relativistic Euler variables
$D=nW$, $M=(\rho+p)W^2v_{\rm flow}$ and
$\mathcal E=(\rho+p)W^2-p-K$.
The constant floor is subtracted only from the numerical energy variable.
The code evolves current, momentum and total energy in conservative form,
then reconstructs the thermal energy from~\eqref{eq:warm}. It imposes no
density or thermal-energy clipping.

Across 128, 256, 512, 1024 and 2048 cells, the largest mean energy drift
was $4.9\times10^{-15}$. The proper density stayed positive. Thermal
energy was nonnegative within pressure-inversion roundoff, whose most
negative value was $-9.9\times10^{-12}$. The maximum tested sound speed
squared was $0.983$, and the warm gate stayed in $[0,1]$ within roundoff.
A separate entropy diagnostic increased throughout its sampled history.
For $\gamma=4/3$, one allowed entropy normalization has comoving entropy
density $e_{\rm th}^{3/4}$.

\begin{figure}[ht]
\centering\includegraphics[width=\textwidth]{thermal_interface_refined.png}
\caption{A selected isotropic thermal continuation keeps compressive
energy in the evolved fluid and feeds its pressure back into the motion.
The finest displayed calculation has 2048 cells.}
\end{figure}

This closes the resolved-energy ledger within the chosen fluid model.
It does not yet provide a converged prediction of heat production: the
mean thermal energy at the five resolutions was $0.0641$, $0.0429$,
$0.0289$, $0.0198$ and $0.0140$ in test units. Numerical viscosity remains
material. Its physical replacement and the relation to an anisotropic
kinetic stock must be matched before using this for a void or halo
prediction. The warm spatial test and the particle-contact test are
separate benchmarks of compatible balance laws, not a combined galaxy
simulation.


\section{What remains before applying the contact to galaxies}
The construction supplies a route from a positive orbital target to local
redistribution with protected angular data and finite work. Three physical
choices still control its predictions.

First, the target $h(E,\ell;H)$ must be selected and made
self-consistent with the evolving potential. The existing positive orbit
weights give a feasible starting point, not a unique distribution or an
assembly history. The coupled model has not resolved the high-redshift
$a_0$ tension.

Second, the matrix elements and bath correlation spectrum must set a
common rate and its spatial range. Detailed balance fixes a ratio, while
equation~\eqref{eq:Hint} explicitly leaves $g_{ij}$ free. Choosing a rate
separately for each galaxy would not produce an independent prediction.
The required weak-coupling and short-memory assumptions need checking at
the assumed horizon temperature.

Third, the selected warm fluid stress must be matched to the anisotropic
kinetic distribution and bath. The fluid and kinetic descriptions are
different resolutions of the same stock; adding both full energy densities
would double count it. A combined implementation must identify their
moments, carry the bath impulse and reproduce the conservative interface
benchmark without introducing an independent full halo.

The new contact can therefore be stated as an effective completion
hypothesis with reproducible equations. It should not be described as
already derived from native microstructure, as a fully established
quantum gravitational response, or as a successful new empirical fit.

\section{Verification and preservation}
Thirty-six standalone Lean certificates check the current-rate
factorization, finite-supply bound, population and impulse identities,
protected-noise algebra, event admissibility, joint energy shifts and
symmetric gate bounds, the positive edge completion, and the warm
enthalpy, floor and sound-speed bounds. Their axiom audit
uses only standard logical axioms.
The available Lean/mathlib 4.19 runtime was used; upstream 4.35 was not rebuilt.
The Hamiltonian, kinetic and spatial simulations have separate verification
records. The global covariant variation, quantum bath reduction and
nonlinear PDE existence are not claimed as Lean theorems.

The preceding research, four publication sources and upstream Lean files
were kept unchanged. New sources, failed checks, earlier gates, numerical
results and build versions are retained in the dated research directory.

\begin{thebibliography}{9}
\bibitem{Seifert} U. Seifert, ``Stochastic thermodynamics, fluctuation theorems,
and molecular machines,'' Rep. Prog. Phys. 75, 126001 (2012),
\href{https://arxiv.org/abs/1205.4176}{arXiv:1205.4176}.
\bibitem{Garriga} J. Garriga and V. F. Mukhanov, ``Perturbations in k-inflation,''
Phys. Lett. B 458, 219 (1999),
\href{https://arxiv.org/abs/hep-th/9904176}{arXiv:hep-th/9904176}.
\bibitem{Brown} J. D. Brown, ``Action functionals for relativistic perfect
fluids,'' Class. Quantum Grav. 10, 1579 (1993),
\href{https://arxiv.org/abs/gr-qc/9304026}{arXiv:gr-qc/9304026}.
\end{thebibliography}
\end{document}
